\clearpage
\section{Protocol details, conditional bounds and authored benchmark}
\label{app:audit-details}
\paragraph{Report authentication and its limits.}
A build identifier and MAC let the center check report origin and integrity.
They do not establish that a build observed all actual inputs, produced correct
labels, or executed in isolation. In the software baseline, a build fingerprint
is an identifier rather than hardware evidence of execution.

\begin{proposition}[Report authenticity, not semantic soundness]\label{prop:attest}
Assume an EUF-CMA-secure MAC, an uncompromised issuer key, and a verifier that
checks a canonical report serialization. An adversary without that key cannot
produce an accepted message--tag pair for a new message not previously
MACed by the issuer, except with negligible probability.
\end{proposition}
\begin{proof}
Such a pair constitutes a MAC forgery under the stated security experiment.
If the implementation authenticates a digest rather than the serialization,
the argument additionally requires collision resistance of that digest.
\end{proof}
This statement allows replay of a previously authenticated report; event binding
and deduplication are separate checks. It also allows an authorized issuer to
sign an incorrect claim. In particular, the proposition proves neither label
accuracy nor completeness of observation. Key isolation and trusted execution
can strengthen provenance under additional assumptions, but do not automatically
make the labeler correct or mediate every communication path. Canonical subject
identifiers prevent some accidental inconsistencies, not false subject claims
by a trusted signer. Section~\ref{sec:methodology} separates these obligations.

\subsection{The \ext{} extension}\label{app:ext-detail}
The three roles of the threat model become fields the protocol can carry. Each
privacy-relevant \texttt{Part} carries a label under
\texttt{Part.metadata["a2a.privacy/v1"]}:
\[
\textsf{PrivacyLabel} = \{\,
\text{subject},\ \text{owner},\ \text{sensitivity}\!\in\![0,5],\ \text{category}[],
\]
\[
\text{inferred\_categories}[],\ \text{purpose}[],\ \text{allowed\_recipients}[],\
\text{ttl\_hops},\ \text{provenance\_id}\,\}.
\]
Identifiers are opaque (\texttt{subject:alice}, \texttt{tenant:hospital})---governance
metadata rather than raw message text; their meanings can still reveal sensitive facts. An \texttt{AgentClearance} (an AgentCard declaration)
states the purposes an agent is cleared for, so cross-tenant rules are checkable
from public metadata. Two fields are central to the results:
\textbf{inferred\_categories} is computed \emph{locally} (the local auditor sees
content; ``near the oncology center'' $\rightarrow$ \texttt{health}) and only the
tag travels---the signal the inference detector accumulates without the center
reading content; \textbf{provenance\_id} is a stable datum identity assigned by the
origin and preserved by forwarders even when they re-word the data, so hop/TTL
tracking follows a datum across paraphrasing (a content hash breaks the moment a relay rephrases).
Figure~\ref{fig:label} maps the implemented fields to their policy checks.
% Shared, editable TikZ. Data: codex_pack/figure_data.json.
\begin{widefigure}[t]\centering
\diagramcontent{%
\begin{tikzpicture}[x=1mm,y=1mm]
\node[figtitle, anchor=west] at (0,65) {\texttt{Part.metadata["a2a.privacy/v1"]}};
\draw[draw=figindigo!50, fill=figindigo!3, rounded corners=1mm, line width=.6pt] (0,1) rectangle (59,60);
\node[figmono, anchor=west] at (2,56.0) {\texttt{data\_subject}};
\node[figmono, anchor=west] at (34,56.0) {\texttt{subject:s}};
\draw[figflow] (60,45.2) -- (64,45.2);
\node[figlabel, anchor=west] at (66,45.2) {Disclosure};
\node[figmono, anchor=west] at (2,50.6) {\texttt{owning\_principal}};
\node[figmono, anchor=west] at (34,50.6) {\texttt{tenant:P}};
\node[figmono, anchor=west] at (2,45.2) {\texttt{allowed\_recipients}};
\node[figmono, anchor=west] at (34,45.2) {\texttt{[tenant:P]}};
\node[figmono, anchor=west] at (2,39.8) {\texttt{sensitivity}};
\node[figmono, anchor=west] at (34,39.8) {\texttt{4}};
\node[figmono, anchor=west] at (2,34.4) {\texttt{category}};
\node[figmono, anchor=west] at (34,34.4) {\texttt{[health]}};
\node[figmono, anchor=west] at (2,29.0) {\texttt{purpose}};
\node[figmono, anchor=west] at (34,29.0) {\texttt{[care]}};
\draw[figflow] (60,29.0) -- (64,29.0);
\node[figlabel, anchor=west] at (66,29.0) {Purpose};
\node[figmono, anchor=west] at (2,23.6) {\texttt{ttl\_hops}};
\node[figmono, anchor=west] at (34,23.6) {\texttt{2}};
\draw[figflow] (60,20.9) -- (64,20.9);
\node[figlabel, anchor=west] at (66,20.9) {Hop / TTL};
\node[figmono, anchor=west] at (2,18.2) {\texttt{provenance\_id}};
\node[figmono, anchor=west] at (34,18.2) {\texttt{d-7f3a}};
\node[figmono, anchor=west] at (2,12.8) {\texttt{inferred\_categories}};
\node[figmono, anchor=west] at (34,12.8) {\texttt{[health]}};
\draw[figflow] (60,10.1) -- (64,10.1);
\node[figlabel, anchor=west] at (66,10.1) {Inference};
\node[figmono, anchor=west] at (2,7.4) {\texttt{inference\_lambda}};
\node[figmono, anchor=west] at (34,7.4) {\texttt{\{\} (default)}};
\node[fignote, anchor=west] at (0,-4) {Purpose also uses recipient AgentClearance.};
\node[fignote, anchor=west] at (0,-9) {Stable provenance links reworded messages; tags reveal categories.};
\draw[draw=figindigo!60,line width=.5pt] (59.5,56) -- (60.5,56) -- (60.5,34.4) -- (59.5,34.4);
\draw[draw=figindigo!60,line width=.5pt] (59.5,23.6) -- (60.5,23.6) -- (60.5,18.2) -- (59.5,18.2);
\draw[draw=figindigo!60,line width=.5pt] (59.5,12.8) -- (60.5,12.8) -- (60.5,7.4) -- (59.5,7.4);

\draw[figzone,fill=figamber!4,draw=figamber!50] (97,1) rectangle (178,60);
\node[figtitle,anchor=west] at (101,54) {WHAT THE CENTER CAN OBSERVE};
\node[figlabel,anchor=west] at (101,44) {Reported subject, owner and recipient links};
\node[figlabel,anchor=west] at (101,35) {Declared categories and policy constraints};
\node[figlabel,anchor=west] at (101,26) {Accumulated hops and inference evidence};
\draw[figbound] (101,20) -- (174,20);
\node[fignote,anchor=west] at (101,14) {Raw message text is excluded from the report.};
\node[fignote,anchor=west] at (101,8) {Labels and linkability can still reveal information.};
\end{tikzpicture}}
\caption{A privacy label mapped to its policy checks, using the implementation's field names. Example values illustrate the schema; they are not a measured record. Subject and recipient identity join fragments; stable provenance joins relays even when their content hashes change.}\label{fig:label}
\end{widefigure}



\subsection{The center-blind auditor and detectors}\label{app:auditor-detail}
These labels are exactly what the auditor reads, and the content is exactly what
it does not. The auditor desensitizes each labeled \texttt{Part} into a
center-view edge carrying a content hash, the principals, and the
label---no raw text (Figure~\ref{fig:desens}). It
checks for undeclared raw-content token matches in its output, allowing declared
label values, schema tokens and routing identifiers. This lexical check is a
regression diagnostic, not a complete detector of semantic information leakage.

\begin{lemma}[Conditional representation-capacity bound]\label{lem:blind}
Let $S$ denote already disclosed side information. Suppose the remaining
center-visible representation is exactly $V=(D,L)$, where $D$ has $b$ bits and
$L$ takes values in a fixed finite alphabet $\mathcal{L}$. Then
\[ I(M;V\mid S)\le b+\log_2|\mathcal{L}|. \]
All visible fields not included in $S$ must be included in $(D,L)$ for this
bound to apply; unbounded free-text labels do not meet the finite-alphabet premise.
\end{lemma}
\begin{proof}
$I(M;V\mid S)\le H(V\mid S)\le H(D\mid S)+H(L\mid S)
\le b+\log_2|\mathcal{L}|$. No zero-information assumption about a public
hash is required.
\end{proof}
This elementary capacity bound is not a confidentiality or differential-privacy
guarantee. Hashes expose equality and may permit dictionary attacks on low-entropy
messages; labels can disclose sensitive facts directly. Repeated observations,
routing identifiers and external knowledge require separate analysis. In
particular, the bound does not establish computational non-reconstruction.
% Shared, editable TikZ. Data: codex_pack/figure_data.json.
\begin{widefigure}[t]\centering
\resizebox{\linewidth}{!}{%
\begin{tikzpicture}[x=1mm,y=1mm]
\node[figtitle, anchor=west] at (0,61) {ONE SUBJECT, TWO SCHEDULING FRAGMENTS};
\node[figtitle, anchor=west] at (2,53) {Local messages};
\node[figtitle, anchor=west] at (80,53) {Emitted metadata (selected fields)};
\draw[draw=figrose!50, fill=figrose!3, rounded corners=1mm, line width=.6pt] (0,30) rectangle (70,48);
\node[figlabel, anchor=west] at (3,42) {Standing appointment every Tuesday};
\node[figlabel, anchor=west] at (3,36) {afternoon for the next 8 weeks.};
\draw[draw=figamber!50, fill=figamber!3, rounded corners=1mm, line width=.6pt] (79,30) rectangle (178,48);
\node[figmono, anchor=west] at (82,43) {\texttt{hash: 303a15068941cca9}};
\node[figmono, anchor=west] at (82,37) {\texttt{category: [schedule]; sensitivity: 4}};
\node[figmono,text=figindigo, anchor=west] at (82,31.5) {\texttt{inferred\_categories: [health]}};
\draw[figcross] (71,39) -- (78,39);
\draw[draw=figrose!50, fill=figrose!3, rounded corners=1mm, line width=.6pt] (0,8) rectangle (70,26);
\node[figlabel, anchor=west] at (3,20) {Can only meet somewhere near};
\node[figlabel, anchor=west] at (3,14) {the oncology center.};
\draw[draw=figamber!50, fill=figamber!3, rounded corners=1mm, line width=.6pt] (79,8) rectangle (178,26);
\node[figmono, anchor=west] at (82,21) {\texttt{hash: f1a39ba19d396523}};
\node[figmono, anchor=west] at (82,15) {\texttt{category: [schedule]; sensitivity: 4}};
\node[figmono,text=figindigo, anchor=west] at (82,9.5) {\texttt{inferred\_categories: [health, location]}};
\draw[figcross] (71,17) -- (78,17);
\draw[draw=figemerald!50, fill=figemerald!3, rounded corners=1mm, line width=.6pt] (0,-16) rectangle (178,1);
\node[figlabel, anchor=center] at (7,-7) {\icon{infer}{emerald}{5mm}};
\node[figtitle, anchor=west] at (14,-5) {Shared subject + recipient \quad $\Rightarrow$ \quad aggregate health evidence};
\node[figlabel, anchor=west] at (14,-11) {Illustrated uniform rule: $p_0=0.1,\ \lambda=3,\ k=2 \Rightarrow P(A\mid k)=0.50 > 0.40$.};
\end{tikzpicture}}
\caption{Scheduling fragments and their stored auditor output. Raw messages are shown only in the local view; selected hashes and labels are shown in the report view. Under the illustrated uniform-evidence rule, two health-tagged fragments cross the configured threshold. \S\ref{sec:eval} tests this threshold against real recipients.}\label{fig:desens}
\end{widefigure}


Here ``center-blind'' describes the absence of raw message text at the center,
not information-theoretic secrecy of the metadata graph. Four detectors run
over this graph:
\begin{enumerate}[leftmargin=1.4em]
\item \textbf{Cross-tenant disclosure}: sensitive data ($\geq\tau$) reaches a
principal that neither owns it nor is an allowed recipient.
\item \textbf{Purpose limitation}: data reaches an agent whose clearance purposes
are disjoint from the data's permitted purposes.
\item \textbf{Hop/TTL}: a datum (tracked by \texttt{provenance\_id}) is relayed
more than \texttt{ttl\_hops} times; robust to a forwarder re-wording content.
\item \textbf{Cross-tenant inference}: no single edge is a disclosure, but a
recipient principal accumulates converging \texttt{inferred\_categories}
fragments about one subject, letting it infer a sensitive category it was never
authorized for (Figure~\ref{fig:desens}); we ground this in the bound below.
\end{enumerate}

Figure~\ref{fig:deploy} shows the report-only deployment and its authentication boundary. One property of the deployment matters for the fourth detector: fragments about one subject arrive in separate reports from separate agents, so a center that forgot every edge after answering could only compose what one report happened to bundle. The service therefore keeps a retention-bounded metadata graph per workspace (edges, never content) and runs the detectors over the retained graph plus each new report, reporting each violation once; retention is a policy knob---a longer window composes more and holds more metadata---and the cross-report case is exercised at the HTTP entry point in the test suite.
% Shared, editable TikZ. Data: codex_pack/figure_data.json.
\begin{widefigure}[t]\centering
\resizebox{\linewidth}{!}{%
\begin{tikzpicture}[x=1mm,y=1mm]
\draw[draw=figindigo!50, fill=figindigo!3, rounded corners=1mm, line width=.6pt] (0,30) rectangle (94,53);
\node[figtitle, anchor=west] at (3,49) {Tenant P};
\node[figlabel, anchor=center] at (8,39) {\icon{agent}{ink}{4mm}};
\node[figlabel, anchor=west] at (14,39) {App agents};
\draw[figflow,draw=figrose] (32,39) -- (39,39);
\draw[draw=figindigo!50, fill=figindigo!3, rounded corners=1mm, line width=.6pt] (40,33) rectangle (91,46);
\node[figlabel, anchor=west] at (43,42) {\texttt{PrivacyTagger} + \texttt{A2AAuditor}};
\node[figlabel, anchor=west] at (43,36) {\texttt{A2AAttestor}: report MAC};
\draw[draw=figindigo!50, fill=figindigo!3, rounded corners=1mm, line width=.6pt] (0,4) rectangle (94,27);
\node[figtitle, anchor=west] at (3,23) {Tenant Q};
\node[figlabel, anchor=center] at (8,13) {\icon{agent}{ink}{4mm}};
\node[figlabel, anchor=west] at (14,13) {App agents};
\draw[figflow,draw=figrose] (32,13) -- (39,13);
\draw[draw=figindigo!50, fill=figindigo!3, rounded corners=1mm, line width=.6pt] (40,7) rectangle (91,20);
\node[figlabel, anchor=west] at (43,16) {\texttt{PrivacyTagger} + \texttt{A2AAuditor}};
\node[figlabel, anchor=west] at (43,10) {\texttt{A2AAttestor}: report MAC};
\node[figtitle, anchor=west] at (0,58) {LOCAL PROCESSING};
\node[figtitle, anchor=west] at (114,58) {AUDIT SERVICE};
\draw[figbound] (103,2) -- (103,55);
\node[fignote, anchor=north] at (103,0) {report boundary};
\draw[figcross] (94,39) -- (113,39);
\draw[figcross] (94,13) -- (113,13);
\node[figlabel, anchor=center] at (103,45) {$E,\;\mathrm{MAC}$};
\draw[draw=figemerald!50, fill=figemerald!3, rounded corners=1mm, line width=.6pt] (114,4) rectangle (178,53);
\node[figlabel, anchor=center] at (120,46) {\icon{attest}{emerald}{4mm}};
\node[figtitle, anchor=west] at (125,46) {\texttt{A2AVerifier}};
\node[fignote, anchor=west] at (118,39) {Key + configured fingerprint; edge digest};
\draw[figflow] (146,35) -- (146,31);
\node[figlabel, anchor=west] at (118,27) {Four detectors re-run per submitted report};
\node[fignote, anchor=west] at (118,20) {Disclosure / purpose / TTL / inference};
\draw[figflow] (146,16) -- (146,13);
\node[figlabel, anchor=center] at (120,9) {\icon{blind}{slate}{4mm}};
\node[figlabel, anchor=west] at (125,9) {Violations store (no raw-text field)};
\node[fignote, anchor=west] at (0,-5) {$E$: content hash + privacy label + routing principals. Local backend required for local tagging.};
\node[fignote, anchor=west] at (0,-10) {Single organisation: set owning principals equal; purpose and hop policies still apply.};
\end{tikzpicture}}
\caption{Where the report-only deployment runs. Two tenants label and hash locally and send authenticated metadata reports; the center verifies and re-runs the detectors within each submitted edge set. Reports are authenticated with a MAC under a per-build key; confinement of that key to the pinned build is the assumption Proposition~\ref{prop:attest} rests on, and a TEE is what makes it hold against a hostile tenant.}\label{fig:deploy}
\end{widefigure}




\paragraph{A formal inference-gain bound.} We model what a recipient infers about
a subject's sensitive attribute $A\in\{0,1\}$ from $k$ converging fragments
$f_1,\dots,f_k$, with prior $0<p_0=P(A{=}1)<1$, threshold $0<\delta<1-p_0$, and odds $O_0=p_0/(1-p_0)$.

\begin{assumption}[Conditional independence]\label{as:ci}
Given $A$, the fragments are independent, each with likelihood ratio
$\lambda_i=\tfrac{P(f_i\mid A{=}1)}{P(f_i\mid A{=}0)}\ge1$; write $\lambda$ for a
common lower bound $\lambda_i\ge\lambda>1$.
\end{assumption}

\begin{proposition}\label{prop:k}
Under Assumption~\ref{as:ci} the posterior odds are $O_k=O_0\prod_i\lambda_i$, the
posterior is $P(A\mid f_{1:k})=O_k/(1+O_k)$, and the inference gain
$g(k)=P(A\mid f_{1:k})-p_0$ is strictly increasing in $k$. With uniform $\lambda$,
the detector (fires iff $g(k)\ge\delta$) fires exactly for
$k\ge k^{*}=\big\lceil\log_\lambda(O_\delta/O_0)\big\rceil$, where
$O_\delta=(p_0+\delta)/(1-p_0-\delta)$.
\end{proposition}
\begin{proof}
By Bayes in odds form, $O_k=O_0\prod_i\lambda_i$; since each $\lambda_i>1$, $O_k$
is strictly increasing in $k$, and $x\mapsto x/(1+x)$ is strictly increasing, so
$g$ is strictly increasing. With uniform $\lambda$, $g(k)\ge\delta \iff
P(A\mid f_{1:k})\ge p_0+\delta \iff O_k\ge O_\delta \iff O_0\lambda^{k}\ge O_\delta
\iff k\ge\log_\lambda(O_\delta/O_0)$; integrality gives the ceiling.
\end{proof}
For defaults $(p_0,\lambda,\delta)=(0.1,3,0.3)$, $k^{*}=2$: one hint never fires,
two converging fragments do (posterior $0.50$, gain $0.40$), replacing the
heuristic $1-2^{-k}$ with an explicit assumed-model quantity and a closed-form threshold (Figure~\ref{fig:gain});
deployments tune $(p_0,\lambda,\delta)$ and may set $\lambda_i$ from the tagger's
hint specificity (so one high-$\lambda$ fragment can suffice).
\emph{On the assumption:} real fragments are correlated, so
Assumption~\ref{as:ci} is an approximation --- positively-correlated (redundant)
hints can make the product over-credit evidence. Historical recipient probes
also exhibit detector misses, but do not isolate tagger coverage from
misspecification of the likelihood model. The bound is an assumed-model design tool, not a calibrated or tight
guarantee.

% Shared, editable TikZ. Data: codex_pack/figure_data.json.
\begin{figure}[t]\centering
\resizebox{\linewidth}{!}{%
\begin{tikzpicture}[x=1mm,y=1mm]
\node[figtitle, anchor=west] at (0,57) {ILLUSTRATIVE INFERENCE RULE};
\draw[figgrid] (8,5) -- (80,5);
\node[fignote, anchor=east] at (6,5) {0};
\draw[figgrid] (8,15.0) -- (80,15.0);
\node[fignote, anchor=east] at (6,15.0) {0.25};
\draw[figgrid] (8,25.0) -- (80,25.0);
\node[fignote, anchor=east] at (6,25.0) {0.5};
\draw[figgrid] (8,35.0) -- (80,35.0);
\node[fignote, anchor=east] at (6,35.0) {0.75};
\draw[figgrid] (8,45) -- (80,45);
\node[fignote, anchor=east] at (6,45) {1};
\node[fignote, anchor=north] at (8,2) {0};
\node[fignote, anchor=north] at (26,2) {1};
\node[fignote, anchor=north] at (44,2) {2};
\node[fignote, anchor=north] at (62,2) {3};
\node[fignote, anchor=north] at (80,2) {4};
\node[fignote, anchor=north] at (44,-5) {Number of tagged fragments};
\draw[figthresh] (8,21) -- (80,21);
\node[fignote,text=figrose, anchor=north east] at (80,20) {Threshold $p_0+\delta=0.40$};
\draw[figseries=figindigo] (8,9.0) -- (26,15.0);
\draw[figseries=figindigo] (26,15.0) -- (44,25.0);
\draw[figseries=figindigo] (44,25.0) -- (62,35.0);
\draw[figseries=figindigo] (62,35.0) -- (80,41.0);
\filldraw[figdot=figindigo] (8,9.0) circle (1.5pt);
\filldraw[figdot=figindigo] (26,15.0) circle (1.5pt);
\filldraw[figdot=figindigo] (44,25.0) circle (1.5pt);
\filldraw[figdot=figindigo] (62,35.0) circle (1.5pt);
\filldraw[figdot=figindigo] (80,41.0) circle (1.5pt);
\filldraw[figmark=figamber] (26,25) circle (2pt);
\node[figlabel,text=figamber, anchor=west] at (8,40) {One strong hint: $\lambda=9$};
\draw[figflow,draw=figamber] (27,36) -- (26,28);
\node[figlabel, anchor=west] at (47,11) {Uniform $\lambda=3$: $k^*=2$};
\node[fignote, anchor=west] at (0,-14) {$p_0=0.10,\ \delta=0.30$; assumed likelihood ratios.};
\end{tikzpicture}}
\caption{The inference rule in closed form. With uniform likelihood ratio three, the posterior first exceeds 0.40 after two fragments. Assigning likelihood ratio nine to one hint reaches 0.50 immediately; neither threshold establishes when a real recipient can infer the attribute.}\label{fig:gain}
\end{figure}


\subsection{The A2A-MT benchmark}
A2A-MT supplies authored multi-tenant A2A interactions with labels and hashes
as the auditor's input. The benchmark tests consistency with template-defined
policy labels, rather than independently adjudicated semantic confidentiality.
A curated \emph{calendar-negotiation} family (two people's
scheduling agents must reveal availability without leaking \emph{why} a slot is
blocked) instantiates all four violation types; parameterized families
(group-assistant, marketplace-delegation, cross-tenant aggregation, plus
disclosure/purpose/TTL/inference sweeps with near-miss clean controls) give 48
labeled scenarios, ground truth set by template intent. An LLM fills each Part
with realistic text matching its label (labels remain the ground truth), exported
as a shareable artifact.

WeClawArena~\cite{weclaw} studies cross-user collaboration in an auditable
sandbox. Our authored scenarios instead restrict the central observer to
metadata. These are different observation settings; we have not performed a
matched comparison and neither suite supplies a numerical upper bound for the
other. The three-role schema is a design choice whose utility must be evaluated
under explicit deployment assumptions.
